% Blueprint: the exact Riemann--von Mangoldt formula
%
%   N(T) = theta(T)/pi + 1 + S(T)
%
% relating the zero-counting function N(T) (mathlib/PNT+ `riemannZeta.N`) to the
% Riemann--Siegel theta function and the argument S(T) = (1/pi) arg zeta(1/2+iT).

\chapter{The exact Riemann--von Mangoldt formula}
\label{cha:n_formula}

The goal of this chapter is the exact (not merely asymptotic) Riemann--von Mangoldt
formula
\[
  N(T) \;=\; \frac{\theta(T)}{\pi} \;+\; 1 \;+\; S(T),
  \qquad
  S(T) := \frac{1}{\pi}\,\Delta_{L}\arg\zeta ,
\]
valid whenever $T>0$ is not the ordinate of a zero of $\zeta$. Here $N(T)$ is the
\emph{zero-counting function} \texttt{riemannZeta.N} of
\texttt{PrimeNumberTheoremAnd} (the number of zeros of $\zeta$ with imaginary part
in $(0,T)$, counted with multiplicity), $\theta$ is the Riemann--Siegel theta
function, and $L$ is the polygonal path $2 \to 2+iT \to \tfrac12+iT$.

This identity is the exact counting tool underlying the certified
Euler--Maclaurin zero-finder: combined with sign changes of Hardy's
$Z(t) = e^{i\theta(t)}\zeta(\tfrac12+it)$ (computed in the reference implementation,
\texttt{python/src/euler\_maclaurin/zeta\_zeros.py}) it certifies that all zeros up to height
$T$ have been found, by matching the algorithmic count to $N(T)$.

The argument follows Titchmarsh, \emph{The Theory of the Riemann Zeta-Function},
\S 9.3, and Edwards, \emph{Riemann's Zeta Function}, \S 6.5--6.7.

\section{Prerequisites from mathlib and PrimeNumberTheoremAnd}
\label{sec:nf_prereqs}

The following are used as black boxes. They fix notation and make the dependency
edges of the new results explicit.

\begin{definition}[Completed zeta; mathlib]
\label{def:nf_completedZeta}
\lean{completedRiemannZeta}
\leanok
$\Lambda(s) := \texttt{completedRiemannZeta}(s) = \pi^{-s/2}\,\Gamma(s/2)\,\zeta(s)$
(\texttt{riemannZeta\_def\_of\_ne\_zero}, with $\Gamma_{\mathbb R}(s)=\pi^{-s/2}\Gamma(s/2)$
as $\texttt{Gamma}\mathbb{R}$). It is meromorphic with simple poles only at $s=0,1$.
The entire normalisation $\Lambda_0(s) := \texttt{completedRiemannZeta}_0(s)$ satisfies
$\Lambda(s) = \Lambda_0(s) - 1/s - 1/(1-s)$ (\texttt{completedRiemannZeta\_eq}) and is
entire (\texttt{differentiable\_completedZeta}$_0$).
\end{definition}

\begin{theorem}[Functional equation; mathlib]
\label{thm:nf_completedZeta_one_sub}
\lean{completedRiemannZeta_one_sub}
\leanok
For all $s$, $\Lambda(1-s) = \Lambda(s)$, and likewise $\Lambda_0(1-s)=\Lambda_0(s)$
(\texttt{completedRiemannZeta}$_0$\texttt{\_one\_sub}).
\end{theorem}

\begin{proof}
\leanok
In Mathlib as \texttt{completedRiemannZeta\_one\_sub}.
\end{proof}

\begin{definition}[Zero-counting function; PrimeNumberTheoremAnd]
\label{def:nf_N}
\lean{riemannZeta.N}
\leanok
$N(T) := \texttt{riemannZeta.N}(T)
       = \sum_{\rho}\, \operatorname{ord}_\rho(\zeta)$, the sum over zeros $\rho$ of
$\zeta$ with $\operatorname{Im}\rho \in (0,T)$ (and arbitrary real part) of the order
\texttt{riemannZeta.order} of $\zeta$ at $\rho$. Built from
\texttt{riemannZeta.zeroes\_rect} and \texttt{riemannZeta.zeroes\_sum}; finiteness
of the index set on compacta is \texttt{riemannZeta.zeroes\_on\_Compact\_finite}.
\end{definition}

\begin{theorem}[No zeros for $\operatorname{Re} s \ge 1$; mathlib]
\label{thm:nf_zeta_ne_zero_one_le}
\lean{riemannZeta_ne_zero_of_one_le_re}
\leanok
If $\operatorname{Re} s \ge 1$ then $\zeta(s)\neq 0$. Together with the trivial-zero
classification (\texttt{riemannZeta\_neg\_two\_mul\_nat\_add\_one} and the functional
equation \texttt{riemannZeta\_one\_sub}) this pins every nonreal zero of $\zeta$ to
the open critical strip $0<\operatorname{Re} s<1$.
\end{theorem}

\begin{proof}
\leanok
In Mathlib as \texttt{riemannZeta\_ne\_zero\_of\_one\_le\_re}.
\end{proof}

\begin{theorem}[Rectangle contour scaffolding; PrimeNumberTheoremAnd]
\label{thm:nf_rectangle}
\lean{RectangleIntegral}
\leanok
\texttt{ResidueCalcOnRectangles.lean} provides the boundary contour integral
$\texttt{RectangleIntegral}\,f\,z\,w$ of $f$ over $\partial R$ for the rectangle $R$
with opposite corners $z,w$, its splittings into horizontal/vertical segments
(\texttt{HIntegral}, \texttt{VIntegral}), Cauchy's theorem on a rectangle for
holomorphic integrands (\texttt{HolomorphicOn.vanishesOnRectangle}), and the
pole-extraction lemma \texttt{RectanglePullToNhdOfPole}. These are the building
blocks for the argument principle below.
\end{theorem}

\begin{proof}
\leanok
In \texttt{PrimeNumberTheoremAnd} (\texttt{ResidueCalcOnRectangles.lean}).
\end{proof}

\section{The Riemann \texorpdfstring{$\xi$}{xi} function}
\label{sec:nf_xi}

\begin{definition}[Riemann $\xi$]
\label{def:nf_xi}
\lean{riemannXi}
\leanok
Set
\[
  \xi(s) \;:=\; \tfrac12\, s(s-1)\,\Lambda(s).
\]
This is the entire completion of $\zeta$ whose zeros are exactly the nontrivial
zeros of $\zeta$.
\end{definition}

\begin{lemma}[$\xi$ is entire]
\label{lem:nf_xi_entire}
\lean{differentiable_riemannXi}
\leanok
\uses{def:nf_xi, def:nf_completedZeta}
$\xi$ is differentiable on all of $\mathbb{C}$. Moreover
$\xi(s) = \tfrac12 s(s-1)\Lambda_0(s) + \tfrac12$.
\end{lemma}

\begin{proof}
\leanok
\uses{thm:nf_completedZeta_one_sub}
From $\Lambda = \Lambda_0 - 1/s - 1/(1-s)$ (\Cref{def:nf_completedZeta}) and
$s(s-1)/(1-s) = -s$,
\[
  \tfrac12 s(s-1)\Lambda(s)
  = \tfrac12 s(s-1)\Lambda_0(s) - \tfrac12(s-1) + \tfrac12 s
  = \tfrac12 s(s-1)\Lambda_0(s) + \tfrac12 .
\]
The right-hand side is a product/sum of entire functions
($\Lambda_0$ is entire by \texttt{differentiable\_completedZeta}$_0$), hence entire.
The pole-cancelling identity is verified pointwise for $s\neq 0,1$ and extends to
$s\in\{0,1\}$ by continuity; in particular $\xi(0)=\xi(1)=\tfrac12$.
\end{proof}

\begin{lemma}[Functional equation for $\xi$]
\label{lem:nf_xi_one_sub}
\lean{riemannXi_one_sub}
\leanok
\uses{def:nf_xi}
For all $s$, $\xi(1-s) = \xi(s)$.
\end{lemma}

\begin{proof}
\leanok
\uses{thm:nf_completedZeta_one_sub}
$\xi(1-s) = \tfrac12 (1-s)(-s)\,\Lambda(1-s)
          = \tfrac12 s(s-1)\,\Lambda(s) = \xi(s)$,
using $(1-s)(-s) = s(s-1)$ and \Cref{thm:nf_completedZeta_one_sub}.
\end{proof}

\begin{lemma}[Conjugation symmetry of $\xi$]
\label{lem:nf_xi_conj}
\lean{riemannXi_conj}
\leanok
\uses{def:nf_xi}
For all $s$, $\xi(\overline{s}) = \overline{\xi(s)}$. In particular $\xi$ is
real on the real axis, and $\xi(\tfrac12+it)\in\mathbb{R}$ for all $t\in\mathbb{R}$.
\end{lemma}

\begin{proof}
\leanok
\uses{lem:nf_xi_one_sub, lem:nf_xi_entire}
The completed zeta is conjugation symmetric, $\Lambda(\overline{s})=\overline{\Lambda(s)}$:
on $\operatorname{Re} s>1$ each factor of the Dirichlet-series representation
$\Lambda(s)=\pi^{-s/2}\Gamma(s/2)\sum n^{-s}$ is conjugation symmetric, and the identity
theorem (\texttt{AnalyticOn} uniqueness, as in
\texttt{Zeta/EulerMaclaurin/Continuation.lean}) extends this to all of $\mathbb{C}$
(\texttt{Zeta/RiemannXi.lean}).
Since $s\mapsto \tfrac12 s(s-1)$ has real coefficients it is also conjugation
symmetric, so $\xi(\overline{s})=\overline{\xi(s)}$.

For $s=\tfrac12+it$ we have $\overline{s}=\tfrac12-it = 1-s$, so by
\Cref{lem:nf_xi_one_sub}
$\overline{\xi(s)} = \xi(\overline{s}) = \xi(1-s) = \xi(s)$, i.e. $\xi(\tfrac12+it)$
is real.
\end{proof}

\begin{lemma}[$\zeta$ is negative on the real segment $(0,1)$]
\label{lem:nf_zeta_real_neg}
\lean{riemannZeta_ofReal_neg_of_mem_Ioo}
\leanok
\uses{def:nf_completedZeta}
For real $\sigma\in(0,1)$ the value $\zeta(\sigma)$ is real and $\zeta(\sigma)<0$; in
particular $\zeta(\sigma)\neq0$.
\end{lemma}

\begin{proof}
\leanok
\uses{def:nf_completedZeta}
For $\operatorname{Re} s>0$ the analytic continuation of $\zeta$ admits the
fractional-part representation (Titchmarsh \S2.1)
\[
  \zeta(s) \;=\; \frac{s}{s-1} \;-\; s\int_1^\infty \{x\}\,x^{-s-1}\,dx ,
  \qquad \{x\}:=x-\lfloor x\rfloor\in[0,1) ,
\]
the integral converging absolutely for $\operatorname{Re} s>0$ since the integrand is
$O\!\left(x^{-\operatorname{Re} s-1}\right)$. This is the first-order ($p=1$)
Euler--Maclaurin expansion of $\sum n^{-s}$ and is available in the project from the
continuation machinery of \texttt{Zeta/EulerMaclaurin/Continuation.lean}: starting from
\texttt{riemannZeta\_eq} with $M=0$, $N=1$ (whose remainder term
$\texttt{remainderℂIoi}\,2\,(y\mapsto y^{-s})\,1$ is, by
\texttt{remainderℂIoi\_cpow\_eq}, an explicit integral of the periodised Bernoulli
kernel $\widetilde B_2$ against $x^{-s-2}$), one integration by parts in reverse
collapses the $\tfrac12 N^{-s}$ and $\widetilde B_2$ terms into the single
$\widetilde B_1(x)=\{x\}-\tfrac12$ integral; absorbing the constant
$-\tfrac12$ gives the displayed $\{x\}$-form. In Lean this reduction is packaged as
\texttt{riemannZeta\_eq\_order\_one}
($\zeta(s)=\tfrac1{s-1}+\tfrac12+\texttt{remainderℂIoi}\,1\,(y\mapsto y^{-s})\,1$ on
$\operatorname{Re} s>0$, $s\neq1$): the order-$1$ improper remainder
$\texttt{remainderℂIoi}\,1$ is obtained from the order-$2$ continuation
\texttt{riemannZeta\_eq} ($M=0,N=1$) by one reverse integration by parts
\texttt{remainderℂIoi\_succ\_eq}, and \texttt{remainderℂIoi\_cpow\_eq} exhibits it as the
$\widetilde B_1=\{x\}-\tfrac12$ integral against $x^{-s-1}$.

\emph{Sign for real $\sigma\in(0,1)$.} Both $\sigma$ and $\sigma-1$ are real, so
$\zeta(\sigma)\in\mathbb{R}$ (the integral is real). The first term is
$\dfrac{\sigma}{\sigma-1}=-\dfrac{\sigma}{1-\sigma}<0$ because $0<\sigma<1$. For the
integral $J:=\int_1^\infty\{x\}\,x^{-\sigma-1}\,dx$ one has $0<J$ (the integrand is
nonnegative everywhere and strictly positive on $(1,2)$, where $\{x\}=x-1>0$) and
\[
  J\;\le\;\int_1^\infty x^{-\sigma-1}\,dx\;=\;\frac1\sigma\;<\;\infty
\]
(using $0\le\{x\}<1$). Hence $-\sigma J<0$, and
$\zeta(\sigma)=-\dfrac{\sigma}{1-\sigma}-\sigma J<0$.
\end{proof}

\begin{lemma}[{$\xi$ has no zeros on the real segment $[-1,2]$}]
\label{lem:nf_xi_ne_zero_real}
\lean{riemannXi_ofReal_ne_zero_of_mem_Icc}
\leanok
\uses{def:nf_xi, lem:nf_zeta_real_neg, lem:nf_xi_one_sub, lem:nf_xi_zeros}
For real $\sigma\in[-1,2]$ one has $\xi(\sigma)\neq0$.
\end{lemma}

\begin{proof}
\leanok
\uses{lem:nf_zeta_real_neg, lem:nf_xi_one_sub, lem:nf_xi_zeros, thm:nf_zeta_ne_zero_one_le}
Split $[-1,2]$ into cases.
\begin{itemize}
\item $\sigma\in(0,1)$: since $0<\operatorname{Re}\sigma<1$, \Cref{lem:nf_xi_zeros}
  gives $\xi(\sigma)=0\iff\zeta(\sigma)=0$, and $\zeta(\sigma)\neq0$ by
  \Cref{lem:nf_zeta_real_neg}.
\item $\sigma=0$ and $\sigma=1$: $\xi(0)=\xi(1)=\tfrac12\neq0$
  (\texttt{riemannXi\_zero}, \texttt{riemannXi\_one}).
\item $\sigma\in(1,2]$: here $\operatorname{Re}\sigma>1$, so the prefactor
  $\tfrac12\sigma(\sigma-1)\,\pi^{-\sigma/2}\Gamma(\sigma/2)$ is nonzero
  ($\Gamma(\sigma/2)\neq0$ as $\operatorname{Re}(\sigma/2)>0$) and $\zeta(\sigma)\neq0$
  (\Cref{thm:nf_zeta_ne_zero_one_le}); hence
  $\xi(\sigma)=\tfrac12\sigma(\sigma-1)\,\pi^{-\sigma/2}\Gamma(\sigma/2)\,\zeta(\sigma)
  \neq0$.
\item $\sigma\in[-1,0)$: by \Cref{lem:nf_xi_one_sub},
  $\xi(\sigma)=\xi(1-\sigma)$ with $1-\sigma\in(1,2]$, reducing to the previous case.
\end{itemize}
The case $(0,1)$ rests on the non-vanishing of $\zeta$ on the open real interval
(\Cref{lem:nf_zeta_real_neg}): reflecting into $\operatorname{Re}\ge1$ is unavailable
there, since $1-\sigma\in(0,1)$ again.
\end{proof}

\begin{lemma}[Zeros of $\xi$ are the nontrivial zeros of $\zeta$]
\label{lem:nf_xi_zeros}
\lean{riemannXi_eq_zero_iff}
\leanok
\uses{def:nf_xi, def:nf_completedZeta}
For $0<\operatorname{Re} s<1$ we have $\xi(s)=0 \iff \zeta(s)=0$, and the orders agree:
$\operatorname{ord}_s(\xi)=\operatorname{ord}_s(\zeta)$ (as
\texttt{meromorphicOrderAt}). Moreover $\xi$ has no zeros with
$\operatorname{Re} s\ge 1$ or $\operatorname{Re} s\le 0$ off the real axis; in
particular $\xi(0)=\xi(1)=\tfrac12\neq0$. (Non-vanishing on the real segment
$[-1,2]$ is the separate \Cref{lem:nf_xi_ne_zero_real}, which needs the
real-interval input \Cref{lem:nf_zeta_real_neg}.)
\end{lemma}

\begin{proof}
\leanok
\uses{thm:nf_zeta_ne_zero_one_le, lem:nf_xi_one_sub}
On the open strip the factor $\tfrac12 s(s-1)$ and $\pi^{-s/2}\Gamma(s/2)$ are
nonzero and holomorphic (here $\operatorname{Re}(s/2)>0$, so $\Gamma(s/2)\neq0$ by
\texttt{Complex.Gamma\_ne\_zero}), hence $\xi$ and $\zeta$ have the same zeros with
the same orders there. Outside the strip: for $\operatorname{Re} s\ge1$,
$\zeta(s)\neq0$ (\Cref{thm:nf_zeta_ne_zero_one_le}) and the prefactor is nonzero away
from $s=1$ where $\Lambda$ has a pole that the prefactor cancels to $\xi(1)=\tfrac12$;
on the line $\operatorname{Re} s=0$ with $\operatorname{Im} s\neq0$ there are likewise
no zeros (apply \texttt{riemannZeta\_one\_sub}: $\zeta(1-s)\neq0$ since
$\operatorname{Re}(1-s)=1$, and the elementary prefactors $\Gamma$, $\cos(\pi s/2)$ are
nonzero there); for $\operatorname{Re} s\le0$ off the real axis use $\xi(1-s)=\xi(s)$
(\Cref{lem:nf_xi_one_sub}) to reflect into $\operatorname{Re}\ge1$. The boundary edge
values $\xi(0)=\xi(1)=\tfrac12$ come from \Cref{lem:nf_xi_entire}. Non-vanishing on the
\emph{full} real segment $[-1,2]$ is \Cref{lem:nf_xi_ne_zero_real}; the open interval
$(0,1)$ needs it, since the reflection $s\mapsto1-s$ maps $(0,1)$ to itself.
\end{proof}

\begin{lemma}[$N(T)$ counts zeros of $\xi$ in a box]
\label{lem:nf_N_eq_xi_count}
\lean{riemannZeta.N_eq_tsum_riemannXi_order}
\leanok
\uses{def:nf_N, lem:nf_xi_zeros}
For $T>0$,
\[
  N(T) \;=\; \sum_{\substack{\rho:\ \xi(\rho)=0\\ 0<\operatorname{Im}\rho<T}}
                \operatorname{ord}_\rho(\xi),
\]
the sum over zeros of $\xi$ with imaginary part in $(0,T)$, counted with
multiplicity. All such $\rho$ lie in the open strip $0<\operatorname{Re}\rho<1$.
\end{lemma}

\begin{proof}
\leanok
\uses{thm:nf_zeta_ne_zero_one_le, lem:nf_xi_zeros}
By \Cref{def:nf_N}, $N(T)$ sums $\operatorname{ord}_\rho(\zeta)$ over zeros $\rho$ of
$\zeta$ with $0<\operatorname{Im}\rho<T$. The trivial zeros lie on the real axis
($\operatorname{Im}=0$, excluded), and every remaining zero with positive imaginary
part has $0<\operatorname{Re}\rho<1$: zeros with $\operatorname{Re}\ge1$ are ruled out
by \Cref{thm:nf_zeta_ne_zero_one_le}, while those with $\operatorname{Re}\le0$ are
ruled out, away from the real axis, by the functional equation
\texttt{riemannZeta\_one\_sub} (its elementary factors $2$, $(2\pi)^{-s}$, $\Gamma$,
$\cos(\pi s/2)$ are all non-vanishing there). This localization is the Lean lemma
\texttt{riemannZeta\_re\_mem\_Ioo\_of\_eq\_zero\_of\_im\_ne\_zero}. On the strip
\Cref{lem:nf_xi_zeros} identifies the zero sets and preserves the order, so the two
sums agree (after reindexing the index set of $N$, which is
\texttt{riemannZeta.zeroes\_rect}, along this set equality).
\end{proof}

\section{Continuous argument variation}
\label{sec:nf_argvar}

To state $\theta$ and $S$ rigorously we need the net change of the argument of a
non-vanishing holomorphic function along a path; the naive principal value
\texttt{Complex.arg} jumps and is wrong for $T$ large.

\begin{definition}[Continuous argument variation]
\label{def:nf_argvar}
\lean{argVariation}
\leanok
\uses{def:av_argVariation, lem:av_exists, lem:av_unique}
Let $\gamma:[a,b]\to\mathbb{C}$ be a path and $f:\mathbb{C}\to\mathbb{C}$ continuous
with $f(\gamma(u))\neq0$ for all $u$. A \emph{continuous logarithm of $f$ along
$\gamma$} is a continuous $g:[a,b]\to\mathbb{C}$ with $\exp(g(u))=f(\gamma(u))$.
Such $g$ exists and is unique up to an additive constant in $2\pi i\,\mathbb{Z}$
(\Cref{lem:av_exists} and \Cref{lem:av_unique}, built by lifting along the covering $\exp$;
see \Cref{cha:arg_variation}). Define the argument variation
\[
  \Delta_\gamma\arg f \;:=\; \operatorname{Im} g(b) - \operatorname{Im} g(a),
\]
independent of the chosen $g$ (\Cref{def:av_argVariation}).
\end{definition}

\begin{lemma}[Additivity and the logarithmic-derivative formula]
\label{lem:nf_argvar_props}
\lean{argVariation_trans}
\leanok
\uses{def:nf_argvar, lem:av_trans, lem:av_mul, lem:av_reverse, lem:av_integral}
Argument variation is additive under concatenation of paths and under
multiplication of integrands:
$\Delta_\gamma\arg(fg) = \Delta_\gamma\arg f + \Delta_\gamma\arg g$. If $\gamma$ is
$C^1$ and $f$ holomorphic and non-vanishing on its image, then
\[
  \Delta_\gamma\arg f \;=\; \operatorname{Im}\!\int_\gamma \frac{f'(z)}{f(z)}\,dz .
\]
\end{lemma}

\begin{proof}
\leanok
\uses{lem:av_trans, lem:av_mul, lem:av_reverse, lem:av_integral}
Additivity under concatenation is \Cref{lem:av_trans} and reversal is
\Cref{lem:av_reverse}; multiplicativity in the integrand is \Cref{lem:av_mul}
(the continuous logarithm of a product is the sum of the continuous logarithms).
The logarithmic-derivative integral formula is \Cref{lem:av_integral}: the anchored
continuous logarithm $g$ is a $C^1$ primitive of
$u\mapsto f'(\gamma(u))\gamma'(u)/f(\gamma(u))$ (\Cref{lem:av_hasDerivAt}), so
$g(b)-g(a)=\int_\gamma f'/f$ by the fundamental theorem of calculus; taking
imaginary parts gives the claim (the real part being
$\log|f(\gamma(b))|-\log|f(\gamma(a))|$).
\end{proof}

\section{The Riemann--Siegel theta function and \texorpdfstring{$S(T)$}{S(T)}}
\label{sec:nf_theta_S}

\begin{definition}[Riemann--Siegel theta]
\label{def:nf_theta}
\lean{riemannSiegelTheta}
\leanok
\uses{def:nf_argvar}
For $T\ge0$ let $\gamma_\Gamma(u) = \tfrac14 + i\tfrac{u}{2}$, $u\in[0,T]$, the
vertical segment from $\tfrac14$ to $\tfrac14+iT/2$, on which $\Gamma\neq0$ (real
part $\tfrac14>0$). Define
\[
  \theta(T) \;:=\; \Delta_{\gamma_\Gamma}\arg\Gamma \;-\; \frac{T}{2}\log\pi
  \;=\; \operatorname{Im}\log\Gamma\!\bigl(\tfrac14 + i\tfrac{T}{2}\bigr)
        - \frac{T}{2}\log\pi ,
\]
where $\operatorname{Im}\log\Gamma$ denotes the continuous branch with
$\theta(0)=0$ (i.e. anchored at $\arg\Gamma(\tfrac14)=0$). This matches
\texttt{theta} in \texttt{python/src/euler\_maclaurin/gamma\_values.py}.
\end{definition}

\begin{definition}[The function $S(T)$]
\label{def:nf_S}
\lean{argZeta}
\leanok
\uses{def:nf_argvar}
For $T>0$ not an ordinate of a zero of $\zeta$, let $L_T$ be the polygonal path
$2 \to 2+iT \to \tfrac12+iT$. On $L_T$, $\zeta\neq0$ (the vertical part has
$\operatorname{Re}=2>1$, and on the horizontal part $\operatorname{Im}=T$ is not an
ordinate of a zero). Define
\[
  S(T) \;:=\; \frac{1}{\pi}\,\Delta_{L_T}\arg\zeta .
\]
This is the continuous variation of $\arg\zeta$ from the real value at $s=2$ to
$s=\tfrac12+iT$, the classical $S(T) = \tfrac1\pi\arg\zeta(\tfrac12+iT)$.
\end{definition}

\section{The argument principle for a rectangle}
\label{sec:nf_argument_principle}

This is the central analytic ingredient; its rectangle scaffolding
(\Cref{thm:nf_rectangle}) is in \texttt{PrimeNumberTheoremAnd}.

\begin{lemma}[Residue of the logarithmic derivative]
\label{lem:nf_logderiv_residue}
\lean{meromorphicOrderAt_eq_residue_logDeriv}
\leanok
Let $f$ be holomorphic near $z_0$ with $f\not\equiv0$, of order
$m=\operatorname{ord}_{z_0}(f)\ge0$ (\texttt{meromorphicOrderAt}). Then $f'/f$ has a
simple pole at $z_0$ with residue $m$: writing $f(z)=(z-z_0)^m h(z)$ with
$h(z_0)\neq0$, $f'(z)/f(z) = m/(z-z_0) + h'(z)/h(z)$ and $h'/h$ is holomorphic at
$z_0$.
\end{lemma}

\begin{proof}
\leanok
By \texttt{meromorphicOrderAt\_eq\_int\_iff} (or the analytic factorisation at an
isolated zero), $f(z)=(z-z_0)^m h(z)$ with $h$ holomorphic and $h(z_0)\neq0$ on a
neighbourhood. Logarithmic differentiation gives the stated decomposition, and
$h'/h$ is holomorphic since $h\neq0$ near $z_0$.
\end{proof}

\begin{theorem}[Argument principle on a rectangle]
\label{thm:nf_argument_principle}
\lean{rectangle_zero_count_eq_argVariation}
\leanok
\uses{thm:nf_rectangle, lem:nf_logderiv_residue, lem:nf_argvar_props}
Let $R$ be the closed rectangle with corners $z,w$ and let $f$ be holomorphic on a
neighbourhood of $R$ with $f\neq0$ on the boundary $\partial R$. Then the number of
zeros of $f$ inside $R$, counted with multiplicity, equals
\[
  \frac{1}{2\pi i}\oint_{\partial R}\frac{f'(z)}{f(z)}\,dz
  \;=\; \frac{1}{2\pi}\,\Delta_{\partial R}\arg f ,
\]
the second equality being the net argument variation around the (closed) boundary.
\end{theorem}

\begin{proof}
\leanok
\uses{thm:nf_rectangle, lem:nf_logderiv_residue, lem:nf_argvar_props,
      thm:ap_main, lem:ap_finite_zeros, thm:ap_residue_sum, lem:ap_integral_eq_argvar}
This is \Cref{thm:ap_main} of \Cref{cha:argument_principle}, built on the single-pole rectangle residue calculus
\Cref{thm:nf_rectangle}. The zeros of $f$ in $R$ are finite
(\Cref{lem:ap_finite_zeros}: isolated zeros of a nonzero holomorphic function on a
compact set, the rectangle analogue of \texttt{riemannZeta.zeroes\_on\_Compact\_finite}).
Reducing a finite zero set to the single-zero case by axis-parallel subdivision
(\Cref{thm:ap_residue_sum}) and using \Cref{lem:nf_logderiv_residue} for the residue,
the contour integral of $f'/f$ equals $2\pi i\sum_\rho \operatorname{ord}_\rho(f)$.
The equality with $\tfrac{1}{2\pi}\Delta_{\partial R}\arg f$ is
\Cref{lem:ap_integral_eq_argvar} (an application of \Cref{lem:nf_argvar_props} edge by
edge): $\oint_{\partial R} f'/f = \Delta_{\partial R}\log f$ has zero real part on the
closed loop (single-valued $\log|f|$), so it equals $i\,\Delta_{\partial R}\arg f$.
\end{proof}

\section{Reduction to the right boundary by symmetry}
\label{sec:nf_reduction}

\begin{lemma}[Halving the contour]
\label{lem:nf_halving}
\lean{argVariation_riemannXi_rectangle_eq_two_mul}
\leanok
\uses{thm:nf_argument_principle, lem:nf_xi_one_sub, lem:nf_xi_conj, lem:nf_argvar_props,
      lem:nf_xi_zeros, lem:nf_xi_ne_zero_real}
Fix $T>0$ not an ordinate of a zero of $\zeta$, and let $R$ be the rectangle with
corners $-1$ and $2+iT$. Then $\xi\neq0$ on $\partial R$, and
\[
  \tfrac{1}{2\pi}\,\Delta_{\partial R}\arg\xi
  \;=\; \tfrac1\pi\,\Delta_{L_T}\arg\xi ,
  \qquad L_T:\ 2 \to 2+iT \to \tfrac12+iT .
\]
\end{lemma}

\begin{proof}
\leanok
\uses{lem:nf_xi_one_sub, lem:nf_xi_conj, lem:nf_xi_zeros, lem:nf_xi_ne_zero_real}
Non-vanishing of $\xi$ on $\partial R$: the bottom edge $\operatorname{Im}=0$,
$\operatorname{Re}\in[-1,2]$ carries no zeros of $\xi$ (\Cref{lem:nf_xi_ne_zero_real});
the
top edge $\operatorname{Im}=T$ carries none since $T$ is not an ordinate; the right
edge $\operatorname{Re}=2$ has $\zeta\neq0$; the left edge $\operatorname{Re}=-1$
reflects under $s\mapsto1-\overline{s}$ to $\operatorname{Re}=2$, where $\xi\neq0$.

\emph{The relevant symmetry is the reflection across the line
$\operatorname{Re}=\tfrac12$, namely $\rho(s):=1-\overline{s}$, which sends
$x+iy\mapsto(1-x)+iy$ and hence \textbf{preserves} the rectangle $R=[-1,2]\times[0,T]$.}
(Beware: the functional-equation map $s\mapsto1-s$ negates the imaginary part and does
\emph{not} preserve $R$.) Combining $\xi(1-s)=\xi(s)$ (\Cref{lem:nf_xi_one_sub}) with
$\xi(\overline{s})=\overline{\xi(s)}$ (\Cref{lem:nf_xi_conj}) gives
\[
  \xi(\rho(s)) \;=\; \xi(1-\overline{s}) \;=\; \xi(\overline{s}) \;=\;
  \overline{\xi(s)},
  \qquad\text{so}\qquad \arg\xi(\rho(s)) = -\arg\xi(s).
\]
Traverse $\partial R$ counterclockwise starting at the bottom-edge midpoint
$\tfrac12$, and split it at the two midpoints $\tfrac12$ and $\tfrac12+iT$ (where
$\xi\neq0$: $\xi(\tfrac12)>0$, and $\xi(\tfrac12+iT)\neq0$ as $T$ is not an ordinate)
into
\[
  P:\ \tfrac12\to2\to2+iT\to\tfrac12+iT
  \qquad\text{(right half)},\qquad
  Q:\ \tfrac12+iT\to-1+iT\to-1\to\tfrac12
  \qquad\text{(left half)}.
\]
The reflection $\rho$ maps the point set of $P$ onto that of $Q$, and the left half
$Q$ is exactly $\rho\circ P$ \emph{traversed in reverse}. Therefore
\[
  \Delta_{Q}\arg\xi
  = -\,\Delta_{\rho\circ P}\arg\xi
  = -\,\Delta_{P}\arg(\xi\circ\rho)
  = -\,\Delta_{P}\bigl(-\arg\xi\bigr)
  = \Delta_{P}\arg\xi,
\]
using additivity/reversal of argument variation (\Cref{lem:nf_argvar_props}) for the
first equality and $\arg\xi\circ\rho=-\arg\xi$ for the third. Hence
$\Delta_{\partial R}\arg\xi = \Delta_P\arg\xi + \Delta_Q\arg\xi = 2\,\Delta_P\arg\xi$.
Finally the initial segment $\tfrac12\to2$ of $P$ lies on the real axis, where
$\xi>0$ (\Cref{lem:nf_xi_zeros}) so $\arg\xi\equiv0$ contributes $0$; thus
$\Delta_P\arg\xi = \Delta_{L_T}\arg\xi$. Dividing by $2\pi$ gives the claim.
\end{proof}

\section{Splitting the argument of \texorpdfstring{$\xi$}{xi} along \texorpdfstring{$L_T$}{L_T}}
\label{sec:nf_split}

Write $\xi(s) = \tfrac12 s(s-1)\cdot\pi^{-s/2}\cdot\Gamma(s/2)\cdot\zeta(s)$ and apply
multiplicativity of argument variation (\Cref{lem:nf_argvar_props}) along $L_T$.

\begin{lemma}[Argument of the polynomial factor]
\label{lem:nf_arg_poly}
\lean{argVariation_poly_L}
\leanok
\uses{def:nf_argvar, def:nf_S, lem:av_slitPlane, lem:nf_argvar_props}
Along $L_T:2\to2+iT\to\tfrac12+iT$ (with $T>0$), the factor $p(s)=\tfrac12 s(s-1)$
is non-vanishing and
\[
  \Delta_{L_T}\arg p \;=\; \pi .
\]
\end{lemma}

\begin{proof}
\leanok
\uses{lem:av_slitPlane, lem:nf_argvar_props}
$p(s)=\tfrac12 s(s-1)$ vanishes only at $s=0,1$, both off $L_T$. By
multiplicativity (\Cref{lem:nf_argvar_props}), with the nonzero constant
$\tfrac12$ contributing $0$,
\[
  \Delta_{L_T}\arg p
  = \Delta_{L_T}\arg(s\mapsto s) + \Delta_{L_T}\arg(s\mapsto s-1).
\]
Both factors stay in the slit plane $\mathbb{C}\setminus(-\infty,0]$ along $L_T$:
$s$ has $\operatorname{Re} s\in[\tfrac12,2]>0$ throughout, and $s-1$ has either
$\operatorname{Re}(s-1)\ge0$ (on the vertical part $\operatorname{Re}=1$) or
$\operatorname{Im}(s-1)=T\neq0$ (on the horizontal part), so it never meets the
nonpositive real axis. \Cref{lem:av_slitPlane} therefore evaluates each variation
as a difference of principal arguments. With $\arg(2)=\arg(1)=0$,
\[
  \Delta_{L_T}\arg(s\mapsto s) = \arg(\tfrac12+iT),
  \qquad
  \Delta_{L_T}\arg(s\mapsto s-1) = \arg(-\tfrac12+iT).
\]
Now $-\tfrac12+iT = -\overline{(\tfrac12+iT)}$ and $\operatorname{Im}(\tfrac12+iT)
=T>0$; for any $z$ with $\operatorname{Im} z>0$ one has
$\arg z + \arg(-\overline{z}) = \pi$ (writing $z=re^{i\alpha}$ with
$\alpha\in(0,\pi)$, $-\overline{z}=re^{i(\pi-\alpha)}$ with $\pi-\alpha\in(0,\pi)$,
so both arguments are the principal values $\alpha$ and $\pi-\alpha$). Hence
$\Delta_{L_T}\arg p = \arg(\tfrac12+iT)+\arg(-\tfrac12+iT)=\pi$.
\end{proof}

\begin{lemma}[Closed-rectangle argument variation of a non-vanishing holomorphic
function vanishes]
\label{lem:nf_rect_argvar_holo_zero}
\lean{rectangleArgVariation_eq_zero_of_analyticOnNhd}
\leanok
\uses{thm:nf_argument_principle, lem:nf_argvar_props, thm:nf_rectangle}
Let $U\subseteq\mathbb{C}$ be open with $f$ analytic and non-vanishing on $U$, and let
$z,w$ satisfy $\operatorname{Re} z\le\operatorname{Re} w$,
$\operatorname{Im} z\le\operatorname{Im} w$ and $R:=\mathrm{Rectangle}(z,w)\subseteq U$.
Then the boundary argument variation vanishes:
\[
  \Delta_{\partial R}\arg f \;=\; 0 .
\]
\end{lemma}

\begin{proof}
\leanok
\uses{thm:nf_rectangle, lem:ap_integral_eq_argvar}
As $f$ is analytic and non-vanishing on $U$, its logarithmic derivative
$f'/f=\texttt{logDeriv}\,f$ is analytic on $U\supseteq R$. Cauchy's theorem on a
rectangle (\texttt{HolomorphicOn.vanishesOnRectangle} from
\texttt{ResidueCalcOnRectangles.lean}, \Cref{thm:nf_rectangle}) gives
$\texttt{RectangleIntegral}\,(f'/f)\,z\,w=0$. The identity
\texttt{rectangleIntegral\_logDeriv\_eq\_I\_mul\_rectangleArgVariation}
(\Cref{lem:ap_integral_eq_argvar}: every edge of $\partial R$ avoids the zeros of $f$,
since $f\neq0$ throughout $U$) reads
$\texttt{RectangleIntegral}\,(f'/f)\,z\,w = i\,\Delta_{\partial R}\arg f$. Hence
$i\,\Delta_{\partial R}\arg f=0$ and, cancelling $i\neq0$,
$\Delta_{\partial R}\arg f=0$.
\end{proof}

\begin{lemma}[Argument of the $\Gamma$ and $\pi$ factors give $\theta$]
\label{lem:nf_arg_theta}
\lean{argVariation_GammaPi_L}
\leanok
\uses{def:nf_theta, def:nf_argvar, lem:nf_argvar_props, lem:nf_rect_argvar_holo_zero,
      lem:av_reparam, lem:av_slitPlane}
Along $L_T$, the factors $\pi^{-s/2}$ and $\Gamma(s/2)$ are non-vanishing and
\[
  \Delta_{L_T}\arg\bigl(\pi^{-s/2}\Gamma(s/2)\bigr) \;=\; \theta(T) .
\]
\end{lemma}

\begin{proof}
\leanok
\uses{def:nf_theta, lem:nf_argvar_props, lem:nf_rect_argvar_holo_zero, lem:av_reparam,
      lem:av_slitPlane}
Write $\Gamma_{\mathbb R}(s)=\pi^{-s/2}\Gamma(s/2)$ ($=\texttt{Gammaℝ}$). Along $L_T$ both
$\pi^{-s/2}$ and $\Gamma(s/2)$ are non-vanishing ($\operatorname{Re}(s/2)\ge\tfrac14>0$,
so $\Gamma(s/2)\neq0$), so multiplicativity (\Cref{lem:nf_argvar_props}, applied on each
of the two affine segments of $L_T$) gives
\[
  \Delta_{L_T}\arg\Gamma_{\mathbb R}
   = \Delta_{L_T}\arg\bigl(\pi^{-s/2}\bigr) + \Delta_{L_T}\arg\bigl(\Gamma(s/2)\bigr).
\]

\emph{The $\pi$ factor.} $\pi^{-s/2}=\exp\!\bigl(-\tfrac{s}{2}\log\pi\bigr)$ has the
\emph{global} continuous logarithm $g(s)=-\tfrac{s}{2}\log\pi$, so directly from the
definition of the variation (\Cref{def:nf_argvar}),
$\Delta_{L_T}\arg(\pi^{-s/2})=\operatorname{Im} g(\text{end})-\operatorname{Im}
g(\text{start})=-\tfrac{\operatorname{Im} s}{2}\log\pi$ between the endpoints $s=2$
($\operatorname{Im}=0$) and $s=\tfrac12+iT$ ($\operatorname{Im}=T$), namely
$-\tfrac{T}{2}\log\pi$. (The branch is global, so the variation depends only on the
endpoints.)

\emph{The $\Gamma$ factor.} Using the affine identity
$\tfrac12\,\mathrm{lineSeg}(z_0,z_1)=\mathrm{lineSeg}(z_0/2,z_1/2)$ on each segment, the
image of $L_T$ under $s\mapsto s/2$ is the two-segment path
$P:\ 1\to1+iT/2\to\tfrac14+iT/2$, and $\Delta_{L_T}\arg\Gamma(s/2)=\Delta_P\arg\Gamma$.
The defining vertical segment $\gamma_\Gamma(u)=\tfrac14+iu/2$, $u:0\to T$, reparametrises
($u=Tv$, \Cref{lem:av_reparam}) to $Q:\ \tfrac14\to\tfrac14+iT/2$, so by
\Cref{def:nf_theta}
\[
  \Delta_Q\arg\Gamma \;=\; \theta(T)+\tfrac{T}{2}\log\pi .
\]
It remains to prove $\Delta_P\arg\Gamma=\Delta_Q\arg\Gamma$.

\emph{Path-independence via a rectangle.} Let $R'$ be the rectangle with corners
$z=\tfrac14$ and $w=1+iT/2$, i.e. $R'=[\tfrac14,1]\times[0,\tfrac T2]$. Traversed
counterclockwise from $\tfrac14$, its boundary is the concatenation
\[
  \partial R':\quad
  \underbrace{\tfrac14\to1}_{B}\ ,\quad
  \underbrace{1\to1+iT/2\to\tfrac14+iT/2}_{P\ (\text{right then top})}\ ,\quad
  \underbrace{\tfrac14+iT/2\to\tfrac14}_{\text{reverse of }Q}.
\]
$\Gamma$ is analytic and non-vanishing on the open half-plane $\operatorname{Re}>0$
(\texttt{Complex.differentiableAt\_Gamma} away from the non-positive integers, and
\texttt{Complex.Gamma\_ne\_zero}), an open set containing $\overline{R'}$. Hence
\Cref{lem:nf_rect_argvar_holo_zero} gives $\Delta_{\partial R'}\arg\Gamma=0$, and by
additivity and reversal (\Cref{lem:nf_argvar_props}),
\[
  0=\Delta_{\partial R'}\arg\Gamma
   =\Delta_B\arg\Gamma+\Delta_P\arg\Gamma-\Delta_Q\arg\Gamma .
\]
On the bottom edge $B\subseteq(0,\infty)$ the function $\Gamma$ is real and strictly
positive (\texttt{Complex.Gamma\_ofReal} with \texttt{Real.Gamma\_pos\_of\_pos}), hence
lies in the slit plane with $\arg\Gamma\equiv0$; so $\Delta_B\arg\Gamma=0$
(\Cref{lem:av_slitPlane}). Therefore $\Delta_P\arg\Gamma=\Delta_Q\arg\Gamma$.

\emph{Conclusion.} Summing the two factors,
\[
  \Delta_{L_T}\arg\Gamma_{\mathbb R}
   = -\tfrac{T}{2}\log\pi + \Delta_Q\arg\Gamma
   = -\tfrac{T}{2}\log\pi + \Bigl(\theta(T)+\tfrac{T}{2}\log\pi\Bigr)
   = \theta(T). \qedhere
\]
\end{proof}

\section{The main theorem}
\label{sec:nf_main}

\begin{theorem}[Exact Riemann--von Mangoldt formula]
\label{thm:N_formula}
\lean{riemannZeta.N_eq}
\leanok
\uses{lem:nf_N_eq_xi_count, thm:nf_argument_principle, lem:nf_halving,
      lem:nf_arg_poly, lem:nf_arg_theta, def:nf_S, def:nf_theta}
Let $T>0$ and suppose $T$ is not the ordinate of any zero of $\zeta$ (equivalently,
$\zeta(\sigma+iT)\neq0$ for all $\sigma\in\mathbb{R}$). Then
\[
  N(T) \;=\; \frac{\theta(T)}{\pi} \;+\; 1 \;+\; S(T).
\]
\end{theorem}

\begin{proof}
\leanok
\uses{lem:nf_N_eq_xi_count, thm:nf_argument_principle, lem:nf_halving,
      lem:nf_arg_poly, lem:nf_arg_theta, lem:nf_argvar_props}
Let $R$ be the rectangle with corners $-1$ and $2+iT$. By \Cref{lem:nf_N_eq_xi_count}
$N(T)$ equals the number of zeros of $\xi$ inside $R$ (all interior zeros lie in the
strip $0<\operatorname{Re}<1$ and have $0<\operatorname{Im}<T$; none lie on
$\partial R$ by the hypothesis on $T$). The argument principle
(\Cref{thm:nf_argument_principle}) gives
\[
  N(T) \;=\; \frac{1}{2\pi}\,\Delta_{\partial R}\arg\xi
        \;=\; \frac{1}{\pi}\,\Delta_{L_T}\arg\xi
\]
by the symmetry reduction (\Cref{lem:nf_halving}). Now factor
$\xi=\tfrac12 s(s-1)\cdot\pi^{-s/2}\Gamma(s/2)\cdot\zeta$ and use multiplicativity of
argument variation (\Cref{lem:nf_argvar_props}):
\[
  \Delta_{L_T}\arg\xi
  = \Delta_{L_T}\arg\bigl(\tfrac12 s(s-1)\bigr)
  + \Delta_{L_T}\arg\bigl(\pi^{-s/2}\Gamma(s/2)\bigr)
  + \Delta_{L_T}\arg\zeta
  = \pi + \theta(T) + \pi S(T),
\]
using \Cref{lem:nf_arg_poly}, \Cref{lem:nf_arg_theta}, and the definition of $S(T)$
(\Cref{def:nf_S}). Dividing by $\pi$,
\[
  N(T) = \frac{1}{\pi}\bigl(\pi + \theta(T) + \pi S(T)\bigr)
       = \frac{\theta(T)}{\pi} + 1 + S(T). \qedhere
\]
\end{proof}

\begin{remark}[Scope and the ordinate hypothesis]
\label{rem:nf_scope}
The hypothesis that $T$ is not an ordinate of a zero ensures $\xi\neq0$ on
$\partial R$, so the argument principle applies cleanly. For $T$ at an ordinate the
standard convention sets $N(T)=\tfrac12\bigl(N(T^-)+N(T^+)\bigr)$; since
\texttt{riemannZeta.N} uses the open interval $(0,T)$, a zero exactly at height $T$
is excluded, and the formula then holds with $S(T)$ interpreted as the limit from
below. The continuous-argument primitive (\Cref{def:nf_argvar},
\Cref{lem:nf_argvar_props}) is developed in \Cref{cha:arg_variation}, built on the
covering-map lift of $\exp$, and the rectangle argument principle
(\Cref{thm:nf_argument_principle}) in \Cref{cha:argument_principle}. The chapter thus
rests on the completed-zeta functional equation, the $\xi$ bookkeeping, the
continuous-argument primitive, and the \texttt{ResidueCalcOnRectangles} scaffolding.
\end{remark}
